% Source: https://artofproblemsolving.com/community/c3h2304172_a_number_theory_problem
% Stable identity: aops-topic-2304172-q01
Since $n+1$ does not divide $n$, $f(n)$ exists, and $f(n)\ge2$ because $1$ divides $n$.

\textbf{$f(n)$ is a prime power.} Let $m=f(n)$. If $m=ab$ with coprime integers $a,b\ge2$, then $a,b<m$, so both divide $n$; as they are coprime, $ab=m$ divides $n$, a contradiction. So $m$ has only one prime factor: $m=p^e$ with $p$ prime and $e\ge1$.

\textbf{The length is at most $3$.}
\begin{itemize}
\item If $n$ is odd, then $f(n)=2$: length $1$.
\item If $n$ is even, then $m=f(n)\ge3$ is a prime power. If $m$ is odd, then $f(m)=2$: length $2$. If $m=2^e$, then $e\ge2$, so $1$ and $2$ divide $m$ but $3$ does not: $f(m)=3$ and $f(3)=2$: length $3$.
\end{itemize}
So every $n$ has a length, and it is at most $3$. $\blacksquare$

\textbf{Least $n$ of length $3$.} We need $n$ even with $f(n)$ a power of $2$, at least $4$. For $n=2$ and $n=4$, $f(n)=3$ (length $2$); for $n=6$, $f(6)=4$, $f(4)=3$, $f(3)=2$. So the answer is $\boxed{6}$.
